\label{sec:related}
In practice the procedure changes what is reported. Fix the region $\nu$ where
the symmetry is supposed to hold, the scales of the state variables, the
candidate class and a tolerance $\delta$ on the relative defect. Fit a model
class that is assumed to contain the truth, or certify simultaneously over a
nested family whose largest member is (\Cref{prop:classes}). Compute
$U(\theta)$, certify what falls below $\delta$, and report $U$ itself rather
than a binary verdict, because it is the smallest tolerance the data support
and the honest analogue of the singular value that would otherwise be
reported; report the extrapolation factor $\sqrt{\kappa\kappa'}$ of
\Cref{prop:kappa} so the reader can see how much of the target domain
the data actually visited. The reporting checklist in Methods lists every
choice that gives the tolerance its meaning.

The limitations follow from \Cref{tab:guarantees}. The model class is the
binding assumption: validity requires the truth to lie in the fitted class, in
the largest class of a nested family, or within an assumed distance of it.
Paying for that distance can leave nothing to report: on a non-polynomial flow
with an exact rotational symmetry the narrowest bound that accounted for model
error was $0.27$, against a tolerance of $0.05$ (\SFmodelerror), so in practice the
route is a class or dictionary that contains the truth (\STdictionary). A
fitted model can be more symmetric than the system it models, and no analysis
of that model alone can detect the difference. A lack-of-fit test that fails to
reject does not establish correct specification: \Cref{prop:lof} quantifies
the misspecification it misses, and in \Cref{sec:selection} misspecification
missed with probability one half can produce true defects about $1.7$ times the
tolerance. Exact coverage of the coefficient region needs Gaussian
homoskedastic errors; the sandwich regions that allow heteroskedastic or
dependent errors are valid only asymptotically, and in our designs they
undercovered (White's region $0.77$ to $0.88$ at $N = 400$ under every error
distribution tried, the cluster-robust region $0.73$ to $0.92$ with $48$
trajectories under autocorrelated errors, the Newey--West region worse still
on trajectories of at most $100$ points, where the certificate built on it
covered in as few as $80\%$ of replications), so dependent data need many
independent trajectories (\SFdependence, \STdependence). Noisy states are not covered: least squares then targets a
pseudo-true field, which on the hardening Duffing oscillator at state noise
$\sigma_x = 0.08$ is more symmetric than the system, so that the certificate concentrates
on a spurious symmetry as data accumulate, while on the symmetric Hopf system
state noise hides the symmetry instead; integral or weak formulations of
model discovery \citep{messenger2021weak}, ensemble resampling
\citep{fasel2022ensemble} and explicit errors-in-variables methods
\citep{carroll2006measurement} are the natural routes, and none is part of the
guarantee here. When its coefficient region covers, the certificate is
conservative: in the tightness study its interval for the plug-in directions
contained the true defect in every replication at every nominal level from
$0.5$ to $0.99$, and for a generator fixed in advance at an exact symmetry its
coverage was at least $0.93$ at nominal
$0.80$ and $0.985$ at nominal $0.95$. In the degree-$3$ class its certified tolerance stayed a nearly
constant factor of about $8$ above the Le~Cam lower bound across $N$ and
$\sigma$; in the classes of degree $5$ and $7$, examined at one noise level,
the median factor rose to $10.9$ and $13.2$, and the pointwise bound for a
generator fixed in advance grew almost as much, so the gap is not mainly a
cost of simultaneity; it reflects the model class and the looseness of the
two-point bound
(\Cref{sec:resolution}). We treat connected
symmetry groups whose generators lie in a prescribed finite-dimensional class,
excluding translations, for which the relative defect degenerates, and
moderate problem sizes ($n \le 6$, $Q \le 210$, $P \le 36$). Discrete
symmetries, and symmetries visible only after a nonlinear change of coordinates
\citep{liu2022machine,yang2024latent}, are outside the framework, as they are
outside the Lie-algebraic setting generally. Finally, the benchmarks are
simulated deliberately, because measuring an error rate requires knowing the
truth exactly, which a real system does not provide. Of the two real data sets
we analyse, the lynx--hare records are too small to support any symmetry claim
(\SNrealdata), and the drifter velocities certify a physically expected
rotation invariance to a stated tolerance, but only under assumptions the data
cannot check (a model class, an observation model, errors uncorrelated with the
velocity although they contain the unmodelled wind forcing, independence
between drifters, and an asymptotic region with $31$ clusters); neither replaces a setting in which symmetries are both real
and known.

The work sits among several lines of research. The idea that continuous
symmetries are a nullspace runs through work on point clouds
\citep{cahill2023liepca}, conserved quantities \citep{kaiser2018discovering},
trained networks \citep{moskalev2022liegg,gruver2023lie} and dynamic mode
decomposition \citep{baddoo2023physics}; Otto et al.~\citep{otto2025unified} unify these,
extend them to nonlinear group actions on manifolds and to sections of vector
bundles, and add convex penalties that promote symmetry during training. Their
conclusion names our problem: it will be important ``to study the perturbative
effects of noisy data in algorithms to discover and promote symmetry, with the
goal of understanding the effects of problem dimension, noise level, and
amount of data.'' We answer with confidence statements rather than with a
perturbation analysis, because a perturbation bound alone does not tell a user
what to declare; the Weyl-inequality baseline in our experiments is the
perturbation answer, and it is valid but never certifies. They also ask whether
their framework can ``find elements that are nearly symmetric, and to bound the
degree of asymmetry''; \Cref{def:defect} is such a bound and \symcert{}
supplies its confidence statement, in a setting where approximate symmetry is
widely recognised in modelling \citep{wang2022approximately,romero2022learning}.
An alternative line searches for generators by nonlinear optimisation or
adversarial training \citep{desai2022symmetry,yang2023liegan,yang2024latent,
dehmamy2021automatic,liu2022machine}, and a related line learns
invariances as hyperparameters through augmentation or marginal likelihood
\citep{benton2020learning,immer2022invariance}. These
methods return a point estimate of a generator with no error control, and the
question asked here, when a discovered generator may be declared and at what
tolerance, applies to them too: whenever the Lie derivative is bilinear in the
generator coordinates and the model coefficients the machinery transfers,
provided the coefficient region is itself valid (for learned features, that
requires features fixed independently of the data used for the region; see
`Learned features' in \Cref{sec:guarantee}), and
for a fully nonlinear model the ellipsoid must be replaced by a confidence
region obtained some other way, with the logic of \Cref{prop:simultaneous}
unchanged. Once a symmetry is in hand it is enforced as hard constraints or
convex penalties \citep{otto2025unified}, as physics-informed losses
\citep{raissi2019physics,akhound2023lie}, as data augmentation
\citep{brandstetter2022lie}, or inside equation discovery itself
\citep{yang2024symmetry}, and everything downstream inherits the reliability of
the discovery step. On the statistical side, equivalence testing is standard
wherever the interesting hypothesis is the null
\citep{schuirmann1987comparison,berger1996bioequivalence,wellek2010testing};
our contribution is not the idea but the observation that symmetry discovery is
one of those settings, with a computable construction for this functional.
Rank tests \citep{cragg1997inferring,robin2000tests,kleibergen2006generalized}
attack the same nullspace-dimension question and place symmetry in the null.
The selection of a generator is the issue addressed by post-selection and
split-sample inference \citep{berk2013valid,lee2016exact,meinshausen2009pvalues,
chernozhukov2018double,andrews2024inference}; here simultaneity over the whole
candidate class dissolves it more cheaply than splitting, which we verify
empirically, whereas selection of the model class needs the separate treatment
of \Cref{prop:classes}, in line with the general difficulty of inference after
model selection \citep{leeb2005model}. Because a single event covers all
generators at once, no multiplicity correction of the Benjamini--Hochberg
kind~\citep{benjamini1995controlling} is needed. Finally, the shape of the
problem (many moment conditions, a low-dimensional parameter, and a test of
whether they can all be satisfied) is that of a test of overidentifying
restrictions \citep{hansen1982large}, with the difference that here the
restriction of interest is the one we would like to accept rather than reject.

Symmetry discovery is an instance of a pattern that recurs whenever a
scientific claim takes the form ``this quantity is zero'': a rank, a
conservation law, a conditional independence, a null effect. In each, an
estimate is compared to a tolerance, and the tolerance is asked to carry
statistical weight it cannot bear. The remedy is the same each time: state the
tolerance in interpretable units, and require the data to \emph{exclude} the
alternative instead of merely failing to detect it. Whether that remedy is
cheap depends on the functional; for the Lie derivative of a model that is
linear in its coefficients it costs a sequence of trust-region subproblems, and
its reach is set by how honestly the model class and the observation model can
be specified.
